% year: תשפ"ו
% type: בוחן
% semester: א
% official_solution: yes
% source: exams/בחנים/תשפ״ו, בוחן, 9141.pdf
% part: א
% number: 2
% points: 20
% categories: continuity, func-limits

\begin{question}
יהיו $a,b\in\R$. נגדיר פונקציה על ידי
\[
f(x)=\begin{cases}
\frac{x^2-4x+3}{x-1}, & x<1,\\
a, & x=1,\\
2b\cdot\frac{\ln x}{x-1}, & x>1
\end{cases}
\]
לאילו ערכים של $a,b$ הפונקציה רציפה ב-$1$ ? נמקו.
\end{question}

\begin{hint}
בגבול משמאל פרקו את המונה לגורמים. בגבול מימין הציבו $y=x-1$ והשתמשו בגבול היסודי $\frac{\ln(1+y)}{y}\to1$.
\end{hint}

\begin{solution}
הפונקציה רציפה ב-1 אם ורק אם הגבולות החד-צדדיים ב-1 קיימים ושווים ל-$f(1)=a$.

\textbf{הגבול משמאל.} $x^2-4x+3=(x-1)(x-3)$, ועבור $x\neq1$ מותר לצמצם:
\[
\lim_{x\to1^-}f(x)=\lim_{x\to1^-}\frac{(x-1)(x-3)}{x-1}=\lim_{x\to1^-}(x-3)=1-3=-2 ,
\]
כאשר בשוויון האחרון השתמשנו ברציפות הפולינום $x-3$.

\textbf{הגבול מימין.} לפי אריתמטיקה של גבולות (הקבוע $2b$ יוצא מהגבול), ועם ההצבה $y=x-1$ ($y\to0^+$ כאשר $x\to1^+$):
\[
\lim_{x\to1^+}f(x)=2b\cdot\lim_{x\to1^+}\frac{\ln x}{x-1}=2b\cdot\lim_{y\to0^+}\frac{\ln(1+y)}{y}=2b\cdot1=2b ,
\]
לפי הגבול היסודי $\lim_{y\to0}\frac{\ln(1+y)}{y}=1$ (מדף הנוסחאות: $\lim_{y\to0}\frac{\log_a(1+y)}{y}=\frac{1}{\ln a}$ עם $a=e$).

\textbf{מסקנה.} $f$ רציפה ב-1 אם ורק אם $-2=a=2b$, כלומר
\[
\boxed{a=-2,\quad b=-1.}
\]
\end{solution}
